\documentclass[a4paper,twoside]{article}

\usepackage{morewrites}

\usepackage[svgnames,dvipsnames,x11names]{xcolor}
\usepackage{graphicx}
\usepackage[english]{babel}
\usepackage[colorlinks=true, linkcolor=NavyBlue, urlcolor=NavyBlue, citecolor=NavyBlue]{hyperref}
\usepackage[a4paper]{geometry}
\usepackage{ts-fonts}
\usepackage{tcolorbox}
\usepackage{fvextra}
\usepackage{amsmath}
\usepackage{float}
\usepackage{seqsplit}
\usepackage{ts-code}

\usepackage{lastpage}
\usepackage{refcount}
\makeatletter
\def\ps@plain{
\let\@oddhead\@empty
\let\@evenhead\@empty
\def\@oddfoot{
\hfil
\ifnum\getpagerefnumber{LastPage}>1\relax
\thepage
\fi
\hfil}
\let\@evenfoot\@oddfoot
}
\makeatother

\title{Contributing}
\author{}\date{}
\begin{document}
\maketitle
I welcome contributions from the community to help improve TeXSmith! Whether it's reporting bugs, suggesting new features, or submitting code changes, your input is valuable. Here's how you can contribute:
\begin{description}
\item[{Reporting issues}] If you encounter any bugs or issues while using TeXSmith, please \href{https://github.com/yves-chevallier/texsmith/issues}{report them}.
\item[{Suggesting features}] Have an idea for a new feature or improvement? We'd love to \href{https://github.com/yves-chevallier/texsmith/issues}{hear it}!
\item[{Submitting code changes}] If you'd like to contribute code, please fork the repository, make your changes, and submit a \href{https://github.com/yves-chevallier/texsmith/pulls}{pull request}. Make sure to follow the coding style and include tests for any new functionality.
\item[{Improving documentation}] Help us keep the documentation up-to-date and comprehensive by suggesting edits or additions.
\item[{Documentation priorities}] Check the [Release Notes][releasenotes] \& Compatibility page and open issues to see which doc sections need attention when the engine evolves.
\item[{Develop templates}] Create and share your own \LaTeX{} templates for TeXSmith users to use.
\end{description}
TeXSmith is a newly developed project and is not ready for production use yet, but you can test it out and help us improve it.
\section{Run the tests}
\begin{tscode}[lang=bash, engine=pygments]
\begin{Verbatim}[commandchars=\\\{\},breaklines, breakanywhere, commandchars=\\\{\}]
git\PY{+w}{ }clone\PY{+w}{ }https://github.com/yves\PYZhy{}chevallier/texsmith.git
\PY{n+nb}{cd}\PY{+w}{ }texsmith
uv\PY{+w}{ }sync
uv\PY{+w}{ }run\PY{+w}{ }pytest
\end{Verbatim}
\end{tscode}
\section{Build the documentation locally}
\begin{tscode}[lang=bash, engine=pygments]
\begin{Verbatim}[commandchars=\\\{\},breaklines, breakanywhere, commandchars=\\\{\}]
uv\PY{+w}{ }sync\PY{+w}{ }\PYZhy{}\PYZhy{}group\PY{+w}{ }docs
uv\PY{+w}{ }run\PY{+w}{ }mkdocs\PY{+w}{ }serve
\end{Verbatim}
\end{tscode}
\section{Test CI}
To test the Continuous Integration (CI) using GitHub Actions, we need \tscodeinline{act} installed on your local machine:
\begin{tscode}[lang=bash, engine=pygments]
\begin{Verbatim}[commandchars=\\\{\},breaklines, breakanywhere, commandchars=\\\{\}]
curl\PY{+w}{ }\PYZhy{}s\PY{+w}{ }https://raw.githubusercontent.com/nektos/act/master/install.sh\PY{+w}{ }\PY{p}{|}\PY{+w}{ }sudo\PY{+w}{ }bash
git\PY{+w}{ }clone\PY{+w}{ }https://github.com/yves\PYZhy{}chevallier/texsmith.git
\PY{n+nb}{cd}\PY{+w}{ }texsmith
act\PY{+w}{ }\PYZhy{}j\PY{+w}{ }build
\end{Verbatim}
\end{tscode}
This would require a \textbf{lot} of disk space, as it uses Docker containers to simulate the GitHub Actions environment. Select the \emph{medium} size image when prompted.
\end{document}
